\documentclass[10pt]{beamer}

\usepackage{verbatim}
%\usepackage[latin1]{inputenc}
\usepackage{amssymb,amsmath,amsfonts,url}
\usepackage{graphicx,color}
\usepackage{pgf,pgfarrows,pgfnodes,pgfautomata,pgfheaps,pgfshade}
\usepackage[utf8]{inputenc}
\usepackage[french]{babel}

\usepackage[all]{xy}
\usepackage{amstext}
\usepackage{xspace}
\usepackage{mdwtab}


\mode<article> {
  \usepackage{fullpage}
  \usepackage{hyperref}
  \setjobnamebeamerversion{advancedcourse.beamer}
}





\newcommand{\alertonce}[1]{\addtocounter{beamerpauses}{-1}\alert<+>{#1}}
\newcommand{\alerttwice}[1]{\addtocounter{beamerpauses}{-1}\alert<+-+(1)>{#1}}

%\newcommand{\subsubsubsection}[1]{\noindent\textbf{#1}}

\newcommand{\cfbox}[1]{\begin{center}\fbox{#1}\end{center}}
\newcommand{\cbox}[1]{\begin{center} #1 \end{center}}

%%%%%%%%%%%%%%%%%%%%%%%%%%%%%%%%%%%%%%%%%%%%%%%%%%%%%%%%%%%%%%%%%%%%%%%%%%%%%%%%%%%%%%%%%

\graphicspath{{Graphics/}}
%\input def_0.tex

%\input{StyleSlides.tex}
\mode<presentation> {
    \usepackage{beamerthemeBoadilla}

  \beamertemplatetransparentcovereddynamic
  \beamertemplateballitem
}


\title[IN100 Cours 3]{\textbf{Fondements informatiques I}\\
Cours 3: Structures itératives}

\author[]{Sorina Ionica \texttt{sorina.ionica@uvsq.fr} \\~\\Sandrine Vial  \texttt{sandrine.vial@uvsq.fr} }


\institute{}

\date{}


%\setbeamercolor{normal text}{fg=black}

\begin{document}


\frame{\titlepage}


\frame[containsverbatim]{
\frametitle{Les structures itératives (boucles)}


 \textbf{Structure de contrôle} qui \textbf{répète} certaines lignes de code \textbf{jusqu'à
ce qu'une condition déterminée soit satisfaite}.

\vspace{0.5 cm}

\emph{Exemple :} On veut afficher \textit{Programmer c'est cool!} 100 fois.
Comment doit-on procéder ?

\vspace{0.5 cm}

\textbf{!!} Il serait fastidieux de taper l'instruction 100 fois :
\[ \mbox{100 fois :}    \left\{
    \begin{array}{l}
        \mbox{print("Programmer c'est cool!")}\\
        \mbox{print("Programmer c'est cool!")}\\
        \mbox{...}\\
        \mbox{print("Programmer c'est cool!")}\\
    \end{array}
\right.\]


}


\frame[containsverbatim]{
\frametitle{L'itération}

En programmation, une boucle a plusieurs composantes: 

\begin{itemize}
\item
Une  \textbf{variable de contrôle} :  initialisée à une certaine valeur, elle est le point de départ dans la boucle
%  La notion d'itération est utilisée quand on doit exécuter le même traitement un certain nombre de fois de manière connue à l'avance ou non.
\item
Une \textbf{condition de sortie} : l'arrêt de l'itération est déclenché par une condition sur l'état des
  variables dans le script.
  
 \item Un \textbf{itérateur} : incrémente généralement la variable de contröle petit-à-petit à chaque boucle successive, jusqu'à ce que celui-ci remplisse la condition de sortie 
\end{itemize}



\begin{exampleblock}{Deux structures itératives}
\begin{verbatim}
    for
    while
\end{verbatim}
\end{exampleblock}

}





\frame[containsverbatim]{
\frametitle{L'instruction \texttt{for}}

\begin{itemize}
\item Structure itérative qui permet d'exécuter un bloc
de code un nombre prédéterminé de fois.

\item En Python, on itére sur les objets d'une séquence d'objets (liste, chaîne de caractères). 
\end{itemize}

\begin{alertblock}{Syntaxe}
\begin{verbatim}
for nom_variable in sequence_objets : 
    instruction 1
    instruction 2
    ...
    instruction n
\end{verbatim}
\end{alertblock}

\begin{itemize}
\item
  La \textbf{variable de contrôle} (\emph{nom\_variable})  prend une par une les valeurs dans la séquence d'objets. 
\item
Les instructions \textbf{indentées} dans le corps de la boucle sont
exécutées une fois pour chaque valeur de la variable de contrôle.
\end{itemize}



}


\frame[containsverbatim]{\frametitle{L'instruction \texttt{for} avec \texttt{range()}}

\begin{itemize}
\item  La fonction \textbf{range(\ldots)} retourne une séquence de nombres, commençant par une valeur initiale
  \textbf{\emph{valeur\_initiale}} et se terminant par
  \textbf{\emph{valeur\_de\_fin}-1}.   
\item  La variable de contrôle prend les
  valeurs définies par \textbf{range(\emph{valeur\_initiale,
  valeur\_de\_fin})}.
\item Les valeurs de la fonction \textbf{range()} doivent être de type
\textbf{entier}.
\begin{itemize}
\item Par exemple, \texttt{range(1.5, 8.5)} est  incorrect.
\end{itemize} 
\end{itemize}

\vspace{0.5 cm}

\begin{exampleblock}{Exemple}
\begin{verbatim}
print("Afficher les valeurs dans l'intervalle:")
for i in range(1, 11) :
     print("i =", i)
print("Terminé")
\end{verbatim}
\end{exampleblock}


}


\frame[containsverbatim]{
\frametitle{La fonction  \texttt{range()}}

\begin{alertblock}{Trois variantes pur \texttt{range()}}
\texttt{range(valeur\_initiale, valeur\_de\_fin)}\\
\texttt{range(valeur\_de\_fin)}\\
\texttt{range (valeur\_initiale, valeur\_de\_fin, increment)}
\end{alertblock}

\begin{itemize}
 \item  La fonction \texttt{range(valeur\_de\_fin)} est équivalente à la
fonction \texttt{range(0, valeur\_de\_fin)}.
\item  Dans la fonction 
\texttt{range(valeur\_initiale, valeur\_de\_fin, increment)} le paramètre \emph{increment} est
utilisé comme \emph{pas d'avancement}. 
\begin{itemize} \item Il est facultatif, s'il est omis,
le pas sera pris par défaut égal à 1.
\end{itemize}
\end{itemize}



\begin{exampleblock}{Exemple}
\small{
\begin{verbatim}
print('Afficher les valeurs dans l'intervalle:')
for i in range(3, 9, 2) :
      print("i =", i)
print("Terminé")	    
\end{verbatim}
}
\end{exampleblock}

}



\frame[containsverbatim]{
\frametitle{Variable de contrôle}



\begin{exampleblock}{Exemple}
\begin{verbatim}
for i in range(1, 11, 3) :
     print("Avant addition i =", i)
     i+= 1
     print("Après addition i =", i)
print("Dernière valeur, i =", i)
print("Terminé")
\end{verbatim}
\end{exampleblock}

\vspace{0.5 cm}

\begin{itemize}
\item En modifiant la valeur de \emph{i} dans le corps de la boucle, la valeur
modifiée sera utilisée pour le reste de cette itération. 
\item Cependant,
lorsque cette itération se termine, la valeur précédente de \emph{i} est
restaurée.
\end{itemize}

}


\frame[containsverbatim]{
\frametitle{Exemple}


Que fait le programme suivant? 

\begin{exampleblock}{}
\begin{verbatim}
n = int(input("Entrez un nombre : "))
result = 0
for i in range(1, n +1):
       result += i
print("La somme vaut :", result)
\end{verbatim}
\end{exampleblock}


}


\frame[containsverbatim]{\frametitle{La boucle \texttt{for} : autres variantes}

\begin{exampleblock}{}
\begin{verbatim}
liste_de_nombres = [23, 43, 56, 22, 109, 34]
sum = 0
for nb in liste_de_nombres :
    sum += nb
moyenne=sum/len(liste_de_nombres)   
print("La moyenne vaut", moyenne)
\end{verbatim}
\end{exampleblock}

\vspace{0.5 cm}

\textbf{\`A retenir :}

\begin{itemize}
\item La boucle \textbf{for} est utilisée quand on veut faire un traitement
  sur un ensemble de données regroupées dans une séquence.
\end{itemize}

  
}



\frame[containsverbatim]{
\frametitle{La boucle \texttt{while}}

Utilisée lorsqu'on ne connaît pas à l'avance le nombre de fois que le
traitement sera répété. 


\begin{alertblock}{Syntaxe}
\begin{verbatim}
while condition :
    instruction 1
    instruction 2
       ...
    instruction n
\end{verbatim}		
\end{alertblock}

\vspace{0.5 cm}

\begin{itemize}
\item Une \textbf{expression de contrôle} (\emph{condition}), expression
  booléenne, contrôle l'exécution de la boucle.
   \item Cette expression  est testée avant chaque itération.
\item Une ou plusieurs actions \textbf{indentées} à répèter. Les
  instructions dans le corps de la boucle sont exécutées si l'expression
  de contrôle reste satisfaite (\textbf{True}).
\end{itemize}
}



\frame[containsverbatim]{
\frametitle{La boucle \texttt{while}}


\begin{exampleblock}{Exemple}
\begin{verbatim}
i = 1
while i < 10 :
     print("i =", i)
     i = i + 1
print("Terminé")
\end{verbatim}
\end{exampleblock}


\begin{itemize}
\item
  La variable \emph{i} est initialisée à 1. 
\item
  Puisque la condition (i \textless{} 10) est \textbf{True}, le
  programme entre dans la boucle et affiche la valeur de \emph{i}
  (première itération).
\item
  La valeur de \emph{i} s'incrémente par 1 (ligne 4), jusqu'à 10; tant
  que la condition est \textbf{True}.
\item
  Lorsque \emph{i} vaut 10, et lors de la vérification de la condition
  de contrôle, la condition (i \textless{} 10) vaut \textbf{False} et
  donc la boucle se termine.
\end{itemize}



}




\frame[containsverbatim]{\frametitle{Exemple}

Quel est le résultat du code suivant ?

\begin{exampleblock}{}
\begin{verbatim}
sum = 0
i = 1
while i < 10:
     sum = sum + i
     i = i + 1
print(sum)
\end{verbatim}
\end{exampleblock}

\vspace{0.5 cm}

\begin{itemize}
\item Il faut inclure un code dans la boucle qui permet de modifier
l'état de l'expression de contrôle.
\item Il faut faire attention à l'indentation des instructions sous
la boucle.
\end{itemize}



}


\frame[containsverbatim]{\frametitle{Un programme plus complexe}

On veut écrire un script qui demande à l'utilisateur de 
deviner un nombre aléatoire. \\%Tant que l'utilisateur n'a pas réussi à deviner le nombre, le
%programme lui indique dans quelle direction il faut deviner. 

Ce jeu se poursuit jusqu'à la bonne réponse de l'utilisateur.





\begin{exampleblock}{}
\small{
\begin{verbatim}
# To random library: https://docs.python.org/3/library/random.html
import random
nombre_a_deviner = random.randint(1, 100)
reponse_utilisateur = 0
while not reponse_utilisateur == nombre_a_deviner :
        reponse_utilisateur = int(input("Entrez un nombre:"))
        if reponse_utilisateur > nombre_a_deviner :
                  print("Trop grand!")
        elif reponse_utilisateur < nombre_a_deviner :
                  print("Trop petit!")
        else :
                  print("Bravo")
\end{verbatim}
}
\end{exampleblock}


La fonction \texttt{random.randint(min, max)} renvoie un
entier aléatoire supérieur ou égal à \texttt{min} et inférieur ou égal à
\texttt{max}.

}




\frame[containsverbatim]{
\frametitle{Boucle \texttt{for} ou boucle \texttt{while}}

\begin{alertblock}{}
Toute boucle \texttt{for} peut être écrite sous la forme d'une
boucle \texttt{while}.
\end{alertblock}

Prenons un exemple:

\begin{exampleblock}{Avec \texttt{for}}
\begin{verbatim}
for i in range (1, 6) :
     print('Programmer c'est cool!')
\end{verbatim}
\end{exampleblock}

\begin{exampleblock}{Avec \texttt{while}}
\begin{verbatim}
i = 0
while (i < 5) :
    i += 1
    print('Programmer c'est cool!')
\end{verbatim}
\end{exampleblock}

}

\frame[containsverbatim]{\frametitle{Choisir entre \texttt{for} et  \texttt{while}}

\begin{exampleblock}{Exemple}
\begin{verbatim}
for a in "hello world":
     print(a)
\end{verbatim}
\end{exampleblock}

\vspace{0.5 cm}
Donner la variante avec \texttt{while} de ce bloc. 

}





\frame[containsverbatim]{\frametitle{Boucles imbriquées}

Le  joueur de décider à l'avance combien de fois il souhaite
jouer.


\begin{exampleblock}{}
\begin{verbatim}
import random
nb_jeu = int(input("Combien de fois vous voulez jouer? "))
for i in range(0, nb_jeu) :
       print("Le Jeu démarre!")
       nombre_a_deviner = random.randint(1, 100)
       reponse_utilisateur = 0
       while not reponse_utilisateur == nombre_a_deviner :
           reponse_utilisateur = int(input("Entrez une réponse:"))
           if reponse_utilisateur > nombre_a_deviner :
                        print("Trop grand!")
           elif reponse_utilisateur < nombre_a_deviner :
                         print("Trop petit !")
           else :
                         print("Bravo!")
\end{verbatim}
\end{exampleblock}





}






\frame[containsverbatim]{\frametitle{D'autres outils de contrôle : \texttt{continue} }

\begin{alertblock}{}
  Le mot clé \texttt{continue} permet de passer prématurément à
  l'itération suivante de la boucle.
\end{alertblock}


\vspace{0.5 cm}

\begin{exampleblock}{Exemple}
\begin{verbatim}
for i in range(1, 21) :
    if i % 2 == 0 :
        continue
        print("résultat?")
    print(i, "est impair")
print("Terminé")
\end{verbatim}

\end{exampleblock}

}



\frame[containsverbatim]{\frametitle{D'autres outils de contrôle :  \texttt{break}}

 \begin{alertblock}{}
  Le mot clé \texttt{break} permet de \emph{casser} (arrêter)
  l'exécution d'une boucle et de passer à l'instruction suivante.
\end{alertblock}

\vspace{0.5 cm}


\begin{exampleblock}{Exemple}
\begin{verbatim}
for i in range(1, 21) :
    if i % 2 == 0 :
        break
    print(i, "est impair")
print("Terminé")
\end{verbatim}
\end{exampleblock}



}





\end{document}

\frame[containsverbatim]{\frametitle{Pour aller plus loin ...}

\begin{alertblock}
\begin{verbatim}
meteo="pluie"
match meteo: 
	case "pluie": 
		print("parapluie")
	case "soleil"; 
		print("t-shirt")
	case "froid"
		print("manteau")			
\end{verbatim}
\end{alertblock}


